\subsection{Dinamica}
\begin{table}
	\centering
	\caption{Parametros del dron}
	\label{tab:param_dron}
	\begin{tabular}{l c l}
		$m$ & 1.6 & kg \\
		$I_{xx}$ & 0.021915027 & kgm$^2$ \\
		$I_{yy} $ & 0.022124548 & kgm$^2$ \\
		$I_{zz}$ & 0.042349865 & kgm$^2$ \\
		$L$ & 0.23 & m \\
		$K_F$ & 6.11 $\times 10^{-8}$ & N/ms
	\end{tabular}
\end{table}

\subsubsection{Linealización}
Mediante la serie de Taylor \ref{eq:taylor} el modelo del cuadricóptero se puede expresar de forma lineal de la siguiente forma
\begin{equation}
	\label{eq:modelo_lineal}
\begin{bmatrix}
\dot{\mathbf{r}} \\ \ddot{\mathbf{r}} \\ \dot{\boldsymbol{\eta}} \\ \ddot{\boldsymbol{\eta}}
\end{bmatrix}
=
\begin{bmatrix}
\mathit{0}_{3\times3} & I_{3\times3} & \mathit{0}_{3\times3} & \mathit{0}_{3\times3} \\
\mathit{0}_{3\times3} & \mathit{0}_{3\times3} & G(\psi_d) & \mathit{0}_{3\times3} \\
\mathit{0}_{3\times3} & \mathit{0}_{3\times3} & \mathit{0}_{3\times3} & I_{3\times3} \\
\mathit{0}_{3\times3} & \mathit{0}_{3\times3} & \mathit{0}_{3\times3} & \mathit{0}_{3\times3}
\end{bmatrix}
\begin{bmatrix}
\mathbf{r} \\ \dot{\mathbf{r}} \\ \boldsymbol{\eta} \\ \boldsymbol{\dot{\eta}}
\end{bmatrix}
+
\begin{bmatrix}
\mathit{0}_{5\times1} & \mathit{0}_{5\times3} \\
m^{-1} & \mathit{0}_{1\times3}\\
\mathit{0}_{3\times1} & \mathit{0}_{3\times3} \\
\mathit{0}_{3\times1} & I_d
\end{bmatrix}
\begin{bmatrix}
F \\ \boldsymbol{\tau}
\end{bmatrix}
\end{equation}

donde:
\[
\mathbf{r} = \begin{bmatrix}x \\ y \\ z\end{bmatrix}, \quad
\boldsymbol{\eta } = \begin{bmatrix}\phi \\ \theta \\ \psi\end{bmatrix}, \quad
\boldsymbol{\tau} = \begin{bmatrix}\tau_x \\ \tau_y \\ \tau_z\end{bmatrix}
\]
\[
\mathit{G}(\psi_d) = \begin{bmatrix}
0 & g & 0\\
-g & 0 & 0\\
0 & 0 & 0
\end{bmatrix}, \quad
I_d = \operatorname{diag}\left(I_{xx}^{-1}, I_{yy}^{-1}, I_{zz}^{-1}\right)
\]
\subsection{Algoritmos de control}
\subsubsection{PD}
El controlador se ajusta mendiate las funciones de trasferencia, que se optienen mediante \ref{eq:ss2tf}, con
\[
C = \begin{bmatrix}
	I_{3\times3} & \mathit{0}_{3\times3} & \mathit{0}_{3\times3} & \mathit{0}_{3\times3} \\
    \mathit{0}_{3\times3} & \mathit{0}_{3\times3} & I_{3\times3} & \mathit{0}_{3\times3}
\end{bmatrix}
\]
\[
D = \mathit{0}_{6\times4}
\]
\begin{gather}
    G_x(s) = \frac{g}{I_{yy} s^4} \\
    G_y(s) = -\frac{g}{I_{xx} s^4} \\
    G_z(s) = \frac{1}{m s^2} \\
    G_{\phi}(s) = \frac{1}{I_{xx} s^2} \\
    G_{\theta}(s) = \frac{1}{I_{yy} s^2} \\
    G_{\psi}(s) = \frac{1}{I_{zz} s^2}
\end{gather}


Dado que no se puede medir los desplazamientos en $x$ y $y$, solo se realizan 4 de los 6 controladores, control de orientacion y altura.

El control de orientacion se realiza mediante

\subsubsection{PID}

\subsubsection{LQI}
Para realizar seguimiento de una referencia se diseña un controlador de modelo interno \ref{eq:c_mi} y se calculan las ganancias como LQR, este al contener un integrador se convierte en LQI, se calculan las ganacias para los siguietes valores
\[\mathbf{AA} =
\left[\begin{array}{cccc:cccccccc}
	0&0&0&0&1&0&0&0&0&0&0&0\\
	0&0&0&0&0&0&1&0&0&0&0&0\\
	0&0&0&0&0&0&0&1&0&0&0&0\\
	0&0&0&0&0&0&0&0&1&0&0&0\\
	\hdashline
	0&0&0&0&0&1&0&0&0&0&0&0\\
	0&0&0&0&0&0&0&0&0&0&0&0\\
	0&0&0&0&0&0&0&0&0&1&0&0\\
	0&0&0&0&0&0&0&0&0&0&1&0\\
	0&0&0&0&0&0&0&0&0&0&0&1\\
	0&0&0&0&0&0&0&0&0&0&0&0\\
	0&0&0&0&0&0&0&0&0&0&0&0\\
	0&0&0&0&0&0&0&0&0&0&0&0
\end{array}\right]\]

\[\mathbf{BB} =
\left[\begin{array}{cccc}
	0&0&0&0\\
	0&0&0&0\\
	0&0&0&0\\
	0&0&0&0\\
	\hdashline
	0&0&0&0\\
	0.625&0&0&0\\
	0&0&0&0\\
	0&0&0&0\\
	0&0&0&0\\
	0&0.021915027&0&0\\
	0&0&0.022124548&0\\
	0&0&0&0.042349865
\end{array}\right]\]
\[\mathbf{Q} = \operatorname{diag}\left(8, 50, 500, 500, 0.5, 0.5, 200, 200, 200, 10, 10, 10\right)\]
\[\mathbf{R} = \operatorname{diag}\left(0.0001, 0.0001, 0.0001, 0.0001\right)\]

Las ganancias del controlador son:
\[\mathbf{K}\]
\subsection{Resultados.}


Describa los resultados obtenidos. 
Valores de Fuerzas, Cargas, 
Valores de esfuerzos hallados, 
si es el caso las dimensiones calculadas, espesores, longitudes, y 
planos de diseño resultantes.